\chapter{Credit Derivatives}
\label{ch:credit_derivatives}
\chaptermark{Credit derivatives}

Credit derivatives let you trade default risk without holding the
bond.  The most important instrument is the \emph{credit default
swap} (CDS).  This chapter derives CDS pricing from first principles,
shows how a term structure of CDS quotes is \emph{stripped} into a
curve of hazard rates (the exact mirror of the zero-curve bootstrap
in \cref{ch:fixed_income_foundations}), introduces the CDX and
iTraxx index families, and ends with the 2008 Gaussian-copula
failure: the model that priced every senior CDO tranche on Wall
Street and was catastrophically wrong when it mattered.

\begin{backgroundbox}[Prerequisites]
\begin{itemize}
  \item \Cref{ch:options_payoffs}: risk-neutral expectation,
    payoff-then-discount, $D(t) = e^{-\rrf t}$.
  \item \Cref{ch:kelly_risk}: expected loss
    $= \Prob(\text{loss}) \times \E[\text{loss}\mid\text{default}]$,
    and recovery reducing LGD.
  \item \Cref{ch:fixed_income_foundations}: zero-curve bootstrap;
    CDS stripping is the credit analogue.
\end{itemize}
\end{backgroundbox}


%% ================================================================
\section{Default probability and recovery rate}
\label{sec:pd_recovery}
%% ================================================================

A \emph{default} is a corporate (or sovereign) failing to meet a
contractual obligation: a missed coupon, missed principal, or
bankruptcy filing.  Default is the single event that makes a bond
risky.  Two numbers summarise that risk.

\begin{definition}[Default probability]
\index{default probability (PD)|textbf}\index{PD (probability of default)|textbf}
The \emph{cumulative default probability} to horizon $T$, written
$\mathrm{PD}(T)$, is the probability that the reference entity
defaults at some time $\tau \le T$.  The \emph{survival probability}
is the complement,
\[
  S(T) \defeq \Prob(\tau > T) = 1 - \mathrm{PD}(T).
\]
\end{definition}

\begin{definition}[Recovery rate]
\index{recovery rate|textbf}
When a default occurs on a bond with face value $1$, the holder
recovers a fraction $R \in [0,1]$ of face value through the
bankruptcy or workout process.  The \emph{loss given default} is
$\mathrm{LGD} = 1 - R$.
\end{definition}

Two PD flavours matter in \cref{sec:cds}:

\begin{itemize}
  \item \emph{Physical} (or \emph{actuarial}) PD,
    $\Prob(\tau \le T)$: the fraction of similar firms that
    default.  Estimated by rating agencies (Moody's, S\&P, Fitch)
    from historical default data by rating bucket.
  \item \emph{Risk-neutral} PD, $\Qm(\tau \le T)$: implied by
    traded prices under the pricing measure.  Typically larger than
    physical PD by a \emph{credit risk premium} \citeHull.
\end{itemize}

CDS pricing lives entirely in the risk-neutral world.  We drop the
$\Qm$ superscript and write $\mathrm{PD}(T)$, $S(T)$; every default
probability in this chapter is risk-neutral.

\subsection*{Typical recovery rates}

Recoveries depend on seniority.  Rating-agency databases give, for
corporate bonds \citeHull:

\begin{center}
\begin{tabular}{@{}lc@{}}
\toprule
\textbf{Seniority} & \textbf{Typical $R$} \\
\midrule
Senior secured           & 50--60\% \\
Senior unsecured         & 35--45\% \\
Subordinated             & 20--30\% \\
Junior subordinated      & 10--20\% \\
\bottomrule
\end{tabular}
\end{center}

The standard CDS convention is $R = 40\%$ for senior unsecured;
sovereign CDS typically assumes $R = 25\%$.

\paragraph{Worked example: expected loss.}
Suppose a 5-year cumulative default probability is
$\mathrm{PD}(5) = 20\%$ and the recovery rate is $R = 40\%$.  The
expected loss on \$1 of face value over 5 years is
\[
  \mathrm{EL}
  = \mathrm{PD}(5) \times (1 - R)
  = 0.20 \times 0.60
  = 0.12
  = 12\%.
\]
Equivalently, \$1{,}200{,}000 on a \$10M position.  This identity
becomes the CDS par spread once discounting and default timing are
added in.

\begin{keyfactbox}[PD and survival]
For a reference entity with risk-neutral default time $\tau$:
\begin{align*}
  \mathrm{PD}(T) &= \Prob(\tau \le T) \qquad \text{(cumulative)}\\
  S(T) &= 1 - \mathrm{PD}(T) \qquad \text{(survival)}\\
  \mathrm{LGD} &= 1 - R \qquad \text{(loss given default)}\\
  \mathrm{EL}(T) &= \mathrm{PD}(T) \cdot (1 - R)
    \qquad \text{(crude expected loss per unit notional).}
\end{align*}
``Crude'' because this ignores discounting and default timing; the
next section corrects both.
\end{keyfactbox}

\begin{pitfallbox}[Do not confuse physical and risk-neutral PD]
A rating-agency 5-year PD for a BBB-rated firm is roughly 2\%; the
risk-neutral PD implied by the same firm's 5-year CDS is typically
\textbf{4--5 times higher}.  The ratio is the \emph{credit risk
premium}, compensation for bearing jump-to-default risk.  Pricing
a CDS with physical PD mis-marks the contract by hundreds of bp.
\end{pitfallbox}


%% ================================================================
\section{Credit default swaps}
\label{sec:cds}
%% ================================================================

A \emph{credit default swap} is a bilateral contract functioning as
insurance against default of a \emph{reference entity} (e.g.\ Ford
Motor Co.).  In plain English: the protection buyer pays a periodic
premium (the \emph{CDS spread}) and, if the reference entity
defaults on its bonds before maturity, receives par in exchange for
the defaulted bond, exactly like a fire-insurance policy that
pays out only if the house burns down.
The CDS market turned corporate credit into a tradable
asset class in the late 1990s; notional outstanding peaked at \$60
trillion in 2007 and has since shrunk to roughly \$8 trillion
post-crisis \citeHull.

\subsection{Mechanics}

Two counterparties:
\begin{itemize}
  \item the \emph{protection buyer} pays a periodic premium, called
    the \emph{CDS spread} $s$, quoted in basis points per annum of
    notional;
  \item the \emph{protection seller} receives the premium and is
    contractually obliged to pay
    $(1 - R) \times \mathrm{notional}$ if a \emph{credit event}
    occurs on the reference entity before maturity.
\end{itemize}

\paragraph{Premium leg.}  The buyer pays $s \cdot \delta_i$ per unit
notional on scheduled dates $t_1 < t_2 < \ldots < t_n = T$, where
$\delta_i$ is the year fraction between payment dates (typically
$\delta_i = 0.25$, i.e.\ quarterly payments under the standard
IMM-date convention).  Payment continues only while the entity is
alive.  If default occurs at some $\tau \in (t_{i-1}, t_i)$, the
buyer pays an \emph{accrual} from $t_{i-1}$ to $\tau$ and then
stops.

\paragraph{Protection leg.}  At default $\tau$, the seller pays
$(1-R)$ per unit notional.  Under physical settlement the buyer
delivers defaulted bonds (face value = notional) for cash of
notional, recovering $R$ on the bonds and netting $(1-R)$.  Since
2009 the standard is cash settlement via the ISDA auction, which
determines a single recovery value per credit event.

\paragraph{Credit events.}  Under the ISDA Credit Derivatives
Definitions, triggers are bankruptcy, failure to pay, and (non-US)
restructuring.  An ISDA Determinations Committee rules on
occurrence, then holds a public auction to set $R$.

\begin{intuitionbox}
A CDS is default insurance: a premium for a lump-sum $(1-R)$ payout
if the reference entity is ``totalled''.  The fair premium equals
the expected discounted cost of the payout.  Unlike insurance, you
need not own the reference bond.  CDS separated credit exposure
from bond ownership, which made credit a tradable asset class.
\end{intuitionbox}

\begin{figure}[htbp]
\centering
\begin{tikzpicture}[>=Stealth]
  % Time axis
  \draw[->, thick] (-0.3, 0) -- (9.8, 0) node[right, font=\footnotesize] {time};
  \foreach \x/\lbl in {0/0, 2/$t_1$, 4/$t_2$, 6/$t_3$, 8/$T$}
    \draw (\x,-0.12) -- (\x,0.12) node[below=5pt, font=\footnotesize] {\lbl};

  % Party labels (left margin)
  \node[font=\footnotesize, align=right] at (-0.5, 1.8) {Buyer\\pays};
  \node[font=\footnotesize, align=right] at (-0.5,-1.9) {Seller\\pays};

  % Premium leg: upward arrows at each payment date
  \foreach \x in {2, 4, 6, 8}
    \draw[->, thick, blue] (\x, 0.15) -- (\x, 1.6)
      node[above, font=\tiny, blue]{$s\,\delta_i$};

  % Default arrow: conditional, occurs between t2 and t3 (at x=5)
  \draw[->, thick, red, dashed] (5, -0.15) -- (5, -2.2)
    node[below, font=\footnotesize, red, align=center]{$(1{-}R)\,N$};
  \node[red, font=\footnotesize\bfseries] at (5.4, -0.45) {$\tau$};

  % "if default" label
  \node[red, font=\tiny, align=center] at (5, -1.1) {if default\\at $\tau$};

  % Premium payments stop after default
  \draw[gray, dashed, thin] (5, 0) -- (5, 1.8);
  \node[gray, font=\tiny, align=center] at (7, 1.0) {payments stop\\after default};
\end{tikzpicture}
\caption{CDS cashflow structure. The protection \emph{buyer} pays quarterly premiums
$s\,\delta_i$ per unit notional (blue arrows, upward) at each scheduled date, but only
while the reference entity survives. If default occurs at $\tau$, the protection
\emph{seller} pays $(1-R)N$ (red dashed arrow, downward) and all further premium
payments stop. The par spread $s^*$ is set so these two legs have equal present value
at inception.}
\label{fig:cds_cashflows}
\end{figure}


\subsection{Pricing: the two legs}

Let $s$ denote the running premium (per annum, per unit notional),
and let $D(t)$ denote the risk-free discount factor to time $t$.
Assume for now that default times and the risk-free rate are
independent (the standard simplifying assumption).

\subsubsection*{Premium leg present value}

Ignoring accrual, the premium leg pays $s\,\delta_i$ per unit
notional at $t_i$ conditional on survival to $t_i$.  Its PV is
\begin{equation}
\label{eq:cds_pl}
  \mathrm{PL}(s)
  = s \cdot \underbrace{\sum_{i=1}^n \delta_i \, D(t_i) \, S(t_i)}_{A \;\defeq\; \text{risky annuity}}
  = s \cdot A,
\end{equation}
where $A$ is the \emph{risky annuity factor} or \emph{DV01}: the
present value of a stream of \$1-per-annum quarterly coupons paid
while the name is alive.

\subsubsection*{Protection leg present value}

The protection leg pays $(1-R)$ at the default time $\tau$ if
$\tau \le T$.  Its PV is
\begin{equation}
\label{eq:cds_pp}
  \mathrm{PP}
  = (1 - R) \int_0^T D(t)\, f(t)\, dt
  = (1 - R) \int_0^T D(t)\, d[1 - S(t)],
\end{equation}
where $f(t) = -S'(t)$ is the risk-neutral density of the default
time, and $d[1 - S(t)] = f(t)\,dt$ is the probability of default
falling in $[t, t+dt]$.

\subsubsection*{Par spread}

At inception the CDS is traded at the \emph{par spread} $s^\star$
that makes the contract costless to both sides:
$\mathrm{PL}(s^\star) = \mathrm{PP}$, so
\begin{equation}
\label{eq:cds_par}
\boxed{
  s^\star
  = \frac{\mathrm{PP}}{A}
  = \frac{(1 - R) \int_0^T D(t) \, d[1 - S(t)]}
         {\sum_{i=1}^n \delta_i \, D(t_i) \, S(t_i)}.
}
\end{equation}

\begin{keyfactbox}[CDS par spread formula]
The fair running spread of a CDS is the \emph{expected discounted
loss} divided by the \emph{risky annuity}:
\[
  s^\star \;=\; \frac{(1-R)\,\E_\Qm\!\left[D(\tau)\,\ind_{\{\tau \le T\}}\right]}
                       {\E_\Qm\!\left[\sum_i \delta_i\, D(t_i)\, \ind_{\{\tau > t_i\}}\right]}.
\]
Numerator: expected PV of the payoff.  Denominator: expected PV of
\$1-per-annum coupons while alive.
\end{keyfactbox}

\begin{intuitionbox}
The par spread equals \emph{expected discounted loss per unit
notional, per year}.  For small, flat hazard and discount rates the
back-of-envelope identity is
\[
  s^\star \;\approx\; h \cdot (1 - R),
\]
where $h$ is the risk-neutral hazard rate.  ``A 300 bp CDS at 40\%
recovery implies a hazard of about 500 bp'' (500 $\approx$ 300/0.60)
is standard mental arithmetic on a credit desk.
\end{intuitionbox}


\subsection{End-to-end worked example}
\label{sec:cds_worked}

Concrete numbers; the Python in \cref{lst:cds_pricer} reproduces
every figure.

\paragraph{Trade details.}
\begin{itemize}
  \item Reference entity: a BBB-rated corporate, senior unsecured.
  \item Notional: $N = \$10{,}000{,}000$.
  \item Maturity: $T = 5$ years.
  \item Premium frequency: quarterly ($\delta_i = 0.25$), so
    $n = 20$ payment dates at $t_i = 0.25, 0.50, \ldots, 5.00$.
  \item Risk-free zero rate: flat $\rrf = 4\%$ continuously compounded,
    $D(t) = e^{-0.04 t}$.
  \item Hazard rate: flat $h = 3\%$, giving
    $S(t) = e^{-0.03 t}$.  (We formally introduce the hazard rate in
    \cref{sec:hazard_stripping}.)
  \item Recovery: $R = 40\%$, so $1 - R = 0.60$.
\end{itemize}

\paragraph{Step 1: risky annuity.}
Plugging into \cref{eq:cds_pl},
\[
  A = \sum_{i=1}^{20} 0.25 \cdot e^{-(0.04 + 0.03)\, t_i}
    = 0.25 \sum_{i=1}^{20} e^{-0.07\, 0.25 i}
    = 4.18194.
\]
So \$1/year of quarterly premium until default or $T=5$ has PV
\$4.1819 per unit notional.

\paragraph{Step 2: protection leg.}
Using $S(t) = e^{-h t}$, $d[1 - S(t)] = h e^{-h t} dt$, so
\cref{eq:cds_pp} evaluates in closed form:
\[
  \mathrm{PP}
  = (1 - R)\int_0^T h \, e^{-(\rrf + h)t}\, dt
  = (1 - R) \cdot \frac{h}{\rrf + h}\left(1 - e^{-(\rrf + h) T}\right).
\]
Substituting:
\[
  \mathrm{PP}
  = 0.60 \cdot \frac{0.03}{0.07}\left(1 - e^{-0.07 \cdot 5}\right)
  = 0.60 \cdot 0.42857 \cdot 0.29531
  = 0.075937.
\]
So the protection leg is worth 7.59\% of notional.

\paragraph{Step 3: par spread.}
\[
  s^\star
  = \frac{\mathrm{PP}}{A}
  = \frac{0.075937}{4.18194}
  = 0.018158
  = 181.6 \text{ bp}.
\]

\paragraph{Step 4: dollar cash flows.}
On $N = \$10$M the two legs match at par: $\mathrm{PP} \cdot N =
\$759{,}373$ of expected discounted loss, and a quarterly premium
coupon $s^\star \delta N = \$45{,}395$.

\paragraph{Step 5: sanity check.}  The small-rate, small-hazard
approximation says
\[
  s^\star
  \;\approx\; h (1 - R)
  \;=\; 0.03 \cdot 0.60
  \;=\; 180 \text{ bp}.
\]
The exact answer 181.6 bp sits 1.6 bp above the approximation; the
gap is $\rrf/(\rrf+h)$ discounting at finite tenor.  Sharp enough
for mental P\&L.

\begin{prosperitybox}
Prosperity 3 had no defaulting corporates, but ``expected loss =
probability $\times$ severity'' is identical to how top teams sized
stop-losses on adverse-selection-prone market-making
(\cref{ch:top_team_lessons}): a 20\% chance of a 60\% adverse move
on \$10k gives expected loss $0.2 \cdot 0.6 \cdot 10{,}000 =
\$1{,}200$.  If quoting profit is below that, widen or stop.  CDS
par spread is the same idea with hazard rate and discounting.
\end{prosperitybox}

\lstset{style=pythonstyle}
\begin{lstlisting}[caption={Minimal CDS pricer. Reproduces every
number in \cref{sec:cds_worked}.},label=lst:cds_pricer]
import math

def cds_par_spread(T, r, h, R, freq=4):
    """Par spread of a flat-hazard, flat-rate CDS.
    Premium leg pays quarterly while alive; protection leg pays
    (1-R) at default time.  Ignores accrued premium at default."""
    dt = 1.0 / freq
    times = [(i+1)*dt for i in range(int(freq*T))]
    # risky annuity: sum of delta * D(t) * S(t)
    annuity = sum(dt * math.exp(-(r+h)*t) for t in times)
    # protection leg: (1-R) * h/(r+h) * (1 - exp(-(r+h)T))
    protection = (1-R) * h/(r+h) * (1 - math.exp(-(r+h)*T))
    return protection / annuity, annuity, protection

if __name__ == "__main__":
    s, A, PP = cds_par_spread(T=5, r=0.04, h=0.03, R=0.40)
    print(f"annuity       = {A:.4f}")
    print(f"protection    = {PP:.6f}")
    print(f"par spread    = {s*1e4:.1f} bp")
    # annuity       = 4.1819
    # protection    = 0.075937
    # par spread    = 181.6 bp
\end{lstlisting}

\subsection{Credit Greeks}
\label{sec:credit_greeks}

The three first-order sensitivities every credit trader quotes:

\begin{itemize}
  \item \textbf{CS01} (\emph{credit spread oh-one}): the P\&L
    impact of a 1~bp parallel shift in the reference entity's CDS
    spread curve.  Equivalently, one basis point of the risky
    annuity $A$ times notional: $\mathrm{CS01} \approx A \cdot N
    \cdot 10^{-4}$.  For the worked example,
    $\mathrm{CS01} \approx 4.182 \times 10{,}000{,}000 \times
    10^{-4} = \$4{,}182$ per bp.
  \item \textbf{JTD} (\emph{jump-to-default}): P\&L if the entity
    defaults right now.  For a protection buyer, JTD $= (1-R) \cdot
    N$ minus the current premium-leg MtM; for a fresh at-market
    trade, JTD $\approx (1-R) \cdot N$.  In the worked example,
    $0.60 \times \$10\mathrm{M} = \$6\mathrm{M}$, the single
    largest risk in a CDS book.
  \item \textbf{IR01}: P\&L from a 1~bp parallel shift in the
    risk-free curve.  Small on CDS (dominant risk is spread, not
    rates) but non-zero via $D(t)$ in the annuity.
\end{itemize}

Risk reports aggregate CS01 by rating/sector/maturity and JTD by
name to surface \emph{concentration risk}: \$50M JTD to a single
name books \$50M of loss overnight on that credit event.

\begin{intuitionbox}
CS01 is day-to-day P\&L risk (spreads move a few bp/day); JTD is
catastrophic-event risk (rare but binary).  Risk managers look at
CS01 first, regulators at JTD first.  A well-hedged book has small
CS01 and small \emph{name-level} JTD, but portfolio-level JTD
can still be large if names correlate (the CDO problem of
\cref{sec:cdo_copula}).
\end{intuitionbox}

\subsection{The CDS--bond basis}
\label{sec:cds_bond_basis}

In theory, long bond + long CDS protection on the same name yields
the Treasury rate, so the bond's \emph{asset swap spread} (over
Treasuries, expressed as a CDS-equivalent running premium) equals
the CDS par spread:
\[
  \mathrm{basis} \;\defeq\; s^\star_{\mathrm{CDS}} - s_{\mathrm{asw}}
  \;\stackrel{?}{=}\; 0.
\]
In practice the basis is non-zero.  Post-2008 the IG basis has been
\emph{negative} (bond spread exceeds CDS) by 20--80~bp, for
structural reasons:
\begin{itemize}
  \item \emph{Funding}: a CDS is unfunded; a long bond ties up
    balance sheet.  Dealers charge funding into the bond spread.
  \item \emph{Counterparty risk}: a CDS seller can default on the
    payout (AIG 2008), reducing protection value.
  \item \emph{Short-selling frictions}: shorting a corporate bond
    is expensive; buying CDS is frictionless, skewing hedge demand
    toward CDS and compressing its spread.
  \item \emph{Cheapest-to-deliver optionality}: the buyer delivers
    any eligible bond at default, widening CDS vs single-bond ASW.
\end{itemize}
Relative-value desks run a \emph{basis book} going ``long bond,
long CDS protection'' when the basis is abnormally negative.  Like
every carry trade, it blows up when liquidity vanishes (Mar 2020,
Oct 2008).

\begin{pitfallbox}[Running spread versus upfront]
Post-2009 standard CDS (US ``Big Bang'', EU ``Small Bang'') trade
with a fixed coupon (100~bp IG, 500~bp HY) plus an \emph{upfront}
payment bringing NPV to zero.  A quoted number may be a par spread
\emph{or} an upfront; always ask.  Conversion uses the risky
annuity: $U = (s^\star - c) \cdot A$.  Getting this wrong on a
5Y IG name costs 4--5 bp per bp of spread error.
\end{pitfallbox}


%% ================================================================
\section{Hazard rates and stripping}
\label{sec:hazard_stripping}
%% ================================================================

The market quotes CDS at standard tenors (1, 3, 5, 7, 10Y); the
spread-curve shape implies a time-varying hazard rate.  This section
defines the hazard rate, derives survival, and bootstraps
piecewise-constant hazards from observed spreads, the credit
mirror of the zero-curve bootstrap in
\cref{ch:fixed_income_foundations}.

\subsection{Hazard rate: definition}

\begin{definition}[Hazard rate]
\index{hazard rate|textbf}
The (risk-neutral) \emph{hazard rate} $h(t)$ is the instantaneous
default intensity conditional on survival to $t$:
\[
  h(t) \, dt
  \;\defeq\; \Prob\!\left(\tau \in [t, t+dt] \mid \tau > t\right).
\]
Equivalently, $h(t) = -S'(t)/S(t) = -\frac{d}{dt}\log S(t)$.
\end{definition}

Integrating the ODE $dS/S = -h(t)\,dt$ from $0$ to $T$:
\begin{equation}
\label{eq:survival_integral}
  S(T) = \exp\!\left(-\int_0^T h(u)\, du\right).
\end{equation}
Constant $h$ gives $S(T) = e^{-hT}$ and exponentially distributed
$\tau$ with mean $1/h$.  A 3\% hazard gives expected time-to-default
$1/0.03 \approx 33$ years and 5Y cumulative PD $1 - e^{-0.15} =
13.9\%$.

\begin{intuitionbox}
The hazard rate is to default what the short rate is to the zero
curve: $D(T) = \exp(-\int_0^T r\, du)$ vs $S(T) = \exp(-\int_0^T h\,
du)$.  Every technique from \cref{ch:fixed_income_foundations}
(bootstrap, splines, forwards) transfers verbatim.  Credit traders
call $h(t)$ the ``credit short rate''.
\end{intuitionbox}

\subsection{Stripping: the bootstrap}

Given par spreads $s_1^\star, \ldots, s_m^\star$ at maturities
$T_1 < \ldots < T_m$, we want a survival curve consistent with all
quotes.  The standard parameterisation is piecewise-constant hazard:
\[
  h(t) = h_k \qquad \text{for } t \in (T_{k-1}, T_k],
\]
with $T_0 \defeq 0$.  The unknowns are $h_1, h_2, \ldots, h_m$.

\paragraph{Bootstrap procedure.}
\begin{enumerate}
  \item Set $k = 1$.  The shortest quote $s_1^\star$ determines
    $h_1$: solve the one-equation-in-one-unknown problem
    $s_1^\star \cdot A_1(h_1) = \mathrm{PP}_1(h_1)$ for $h_1$ by
    root-finding (1-D Newton or bisection).  Before $T_1$, only
    $h_1$ enters both sides.
  \item For $k = 2, \ldots, m$: $h_1, \ldots, h_{k-1}$ are now
    known.  Solve
    $s_k^\star \cdot A_k(h_k; h_1, \ldots, h_{k-1})
       = \mathrm{PP}_k(h_k; h_1, \ldots, h_{k-1})$
    for the single unknown $h_k$.
\end{enumerate}

Each step is a scalar root-find.  The survival function at any
interior time $t \in (T_{k-1}, T_k]$ is then
\[
  S(t)
  = \exp\!\left(-\sum_{j=1}^{k-1} h_j (T_j - T_{j-1})
                - h_k (t - T_{k-1})\right).
\]

\paragraph{Reducing the root-find.}  Each step is a 1-D Newton on
\[
  g_k(h_k) \defeq s_k^\star \cdot A_k(h_k) - \mathrm{PP}_k(h_k),
\]
with initial guess $h_k^{(0)} = s_k^\star/(1-R)$; it converges in
3--5 iterations.  Derivatives $dA_k/dh_k$, $d\mathrm{PP}_k/dh_k$
have closed forms on piecewise-constant intervals; otherwise use
numerical derivatives.

\paragraph{Sanity checks on a stripped curve.}
\begin{itemize}
  \item $h_k \ge 0$: a negative hazard usually means a stale or
    fat-finger quote (or an exotic sovereign).
  \item $S$ monotone: $S(T_1) \ge S(T_2)$ for $T_1 \le T_2$,
    equivalent to $h_k \ge 0$.
  \item No implausible jumps at tenor boundaries.  A jump from
    200~bp to 50~bp between 3Y and 5Y either prices a known
    corporate event in that window, or signals piecewise-constant
    is too coarse; switch to piecewise-linear or tension splines.
\end{itemize}

\paragraph{Worked sanity-check.}  1Y/3Y/5Y CDS on a BBB name at
100/150/200~bp with $R = 40\%$.  The approximation $h \approx
s/(1-R)$ gives hazards $\approx$ 167, 250, 333~bp: upward-sloping,
consistent with the credit curve.  The exact bootstrap shifts these
by a few bp and preserves shape.

\begin{keyfactbox}[CDS curve stripping]
Given par spreads at maturities $\{T_k\}$ and piecewise-constant
hazard on $(T_{k-1}, T_k]$, the $\{h_k\}$ are recovered by a 1-D
root-find per step, in increasing maturity order.  Mathematically
identical to the zero-curve bootstrap of
\cref{ch:fixed_income_foundations}: at step $k$, prior quantities
are pinned, so one unknown meets one equation.
\end{keyfactbox}

\begin{pitfallbox}[Spread curve shape does not equal hazard curve shape]
If the 1Y and 5Y par spreads are both 200~bp, the hazard curve is
\emph{not} flat; it's slightly downward-sloping, because a flat
long-maturity spread embeds the average of lower near-term hazards.
The spread-to-hazard mapping is nonlinear; only a bootstrap reads
the hazard curve off a price screen.
\end{pitfallbox}


%% ================================================================
\section{CDX and iTraxx indices}
\label{sec:cdx}
%% ================================================================

Individual-name CDS are \emph{single-name CDS}.  Trading a basket
in one wrapper is cheaper and more liquid, so the market built
standardised credit indices.

\begin{definition}[CDS index]
\index{CDS index|textbf}\index{CDX|textbf}\index{iTraxx|textbf}
A \emph{credit index} is a standardised equally-weighted portfolio
of single-name CDS on a fixed roster of entities, trading as one
instrument.  Buying index protection is economically equivalent to
buying protection pro-rata on each name.
\end{definition}

The major families, all administered by Markit/S\&P:
\begin{itemize}
  \item \textbf{CDX.NA.IG}: 125 North-American IG names,
    reconstituted each March/September on ``index roll'' dates.
    The flagship IG benchmark.
  \item \textbf{CDX.NA.HY}: 100 North-American HY names.
  \item \textbf{iTraxx Europe}: 125 European IG names.
  \item \textbf{iTraxx Crossover}: 75 European sub-IG names.
  \item \textbf{iTraxx Asia/Japan, CDX.EM}: regional offshoots.
\end{itemize}

Index CDS trade with fixed coupons (100~bp IG, 500~bp HY) plus an
upfront bringing NPV to zero, matching single-name convention
post-2009.

\begin{historybox}[title={Birth of the CDS: JP Morgan 1994}]
The first CDS was structured in 1994 by Blythe Masters' JP Morgan
team to hedge a \$4.8bn credit line to Exxon (Valdez-spill
liability), paying the EBRD to take the default risk.  Standardised
ISDA documentation arrived by 2000; notional hit \$6tn by 2004 and
\$60tn by 2007.  CDS made credit a tradable asset class; overused
and misunderstood, the same instrument nearly destroyed the system
in 2008.
\end{historybox}

\paragraph{Index arbitrage and the skew.}  The \emph{intrinsic}
level is the average par spread of the 125 constituents; the
\emph{quoted} index spread deviates from it by the \emph{basis} or
\emph{index skew}.  Tight in calm markets (a few bp), it widens in
stress as investors buy index protection (cheap, liquid) faster
than dealers hedge it on the 125 singles.  A dealer arbitrage buys
index, sells singles pro rata, and holds until convergence.

\begin{gsbox}[title={The index trading desk}]
Sell-side credit desks usually run the index book separately from
single-names: indices trade in \$100M clips with hedge funds and
macro, singles in \$5--10M clips with real money.  The index Strat
owns a basis monitor (intrinsic vs quoted for CDX.IG/HY, iTraxx
Main/Xover) and a PnL explain attributing the day's move to index
level, basis carry, single-name hedges, and roll costs.  Interviews
test why the skew exists, when it blows out, and how to monetise a
5~bp dislocation.
\end{gsbox}

\paragraph{Tranches on indices.}  An index can be carved into
\emph{tranches} of different seniority: the subject of the next
section, and of 2008.


%% ================================================================
\section{CDO tranches and the Gaussian copula}
\label{sec:cdo_copula}
%% ================================================================

A \emph{collateralised debt obligation} (CDO) packages a pool of
credit-risky assets (CDS in a \emph{synthetic CDO}, mortgage loans
in a \emph{cash CDO}) and issues layered claims.  Each claim is a
\emph{tranche} defined by an attachment/detachment point in
portfolio-loss space.  Plain English: pool 125 bonds, then sell
slices in order of who eats losses first.  The \emph{equity}
tranche absorbs the first few percent of defaults and is paid the
highest yield as compensation for bearing first-loss risk; the
\emph{senior} tranche pays the lowest yield but is only touched if
losses are catastrophic.

\subsection{Tranche mechanics}

Let $L(t) \in [0,1]$ be the fractional portfolio loss at time $t$.
A tranche with attachment $a$ and detachment $d$ ($0 \le a < d \le
1$) has loss profile
\[
  L_{[a,d]}(t)
  = \frac{\min(L(t), d) - \min(L(t), a)}{d - a} \in [0, 1].
\]
The tranche absorbs nothing below $a$, 1-for-1 losses between $a$
and $d$, and is wiped out above $d$.

Standard CDX.NA.IG tranches:
\begin{center}
\begin{tabular}{@{}lcc@{}}
\toprule
\textbf{Tranche} & \textbf{Attachment} & \textbf{Detachment} \\
\midrule
Equity          & 0\%  & 3\%   \\
Junior mezz     & 3\%  & 7\%   \\
Senior mezz     & 7\%  & 10\%  \\
Senior          & 10\% & 15\%  \\
Super-senior    & 15\% & 30\%  \\
\bottomrule
\end{tabular}
\end{center}
Equity is first-loss and pays the highest coupon; super-senior is
safest and, pre-2008, rated AAA by S\&P and Moody's.

\subsection{The Gaussian copula model}

Tranche pricing needs the \emph{joint} default distribution across
125 names, not marginals alone.  The industry model
\citep{li2000default}, Li's 2000 JP Morgan note, uses a
\emph{single-factor Gaussian copula}.

\paragraph{Construction.}  For each name $i = 1, \ldots, 125$:
\begin{enumerate}
  \item Extract marginal risk-neutral $F_i(t) = \mathrm{PD}_i(t)$
    from the single-name CDS curve.
  \item Draw $X_i = \sqrt{\rho}\, M + \sqrt{1 - \rho}\, Z_i$, with
    systemic factor $M$, idiosyncratic $Z_i$, all iid
    $\Normal(0,1)$, and $\rho \in [0,1]$ the \emph{copula
    correlation}.
  \item Map to default times via $\tau_i = F_i^{-1}(\CDFN(X_i))$.
    Marginals are preserved; the joint distribution is governed by
    $\rho$.
\end{enumerate}
Monte Carlo over $(M, Z_1, \ldots, Z_{125})$ gives the joint
default distribution; tranche par spreads follow from
\cref{eq:cds_par} with $S(t)$ replaced by $\E[1 - L_{[a,d]}(t)]$.

\paragraph{Implied correlation.}  Inverting for the $\rho$ that
reproduces each tranche's market spread gives an \emph{implied
correlation}, analogous to implied vol.  A true Gaussian copula
would price every tranche with a single $\rho$; in reality implied
$\rho$ varies: the \emph{correlation skew}, credit analogue of
the vol smile (\cref{ch:vol_surface}).

\begin{keyfactbox}[Li's single-factor Gaussian copula]
Correlated default times are generated via
\[
  X_i = \sqrt{\rho}\, M + \sqrt{1 - \rho}\, Z_i, \quad
  M, Z_i \sim \Normal(0,1) \text{ iid}, \quad
  \tau_i = F_i^{-1}(\CDFN(X_i)).
\]
One $\rho$ pins the entire correlation structure; tranche prices
follow by Monte Carlo or semi-analytic integration over $M$.
\end{keyfactbox}

\subsection{Semi-analytic tranche loss: a back-of-envelope}
\label{sec:tranche_semianalytic}

Conditional on $M = m$, the defaults $\{\tau_i \le T\}$ are
independent Bernoulli events with common probability
\[
  p(m)
  \;=\; \Prob\!\left(X_i < \CDFN^{-1}(\mathrm{PD}(T)) \mid M = m\right)
  \;=\; \CDFN\!\left(\frac{\CDFN^{-1}(\mathrm{PD}(T)) - \sqrt{\rho}\,m}{\sqrt{1 - \rho}}\right).
\]
The conditional default count is binomial; for large portfolios
the \emph{large-homogeneous-pool} (LHP) approximation \citeHull\
replaces it by its mean $p(m)$.  Integrating over $M \sim
\Normal(0,1)$ gives the unconditional loss density.

\paragraph{Numeric illustration.}  $\mathrm{PD}(5) = 5\%$, $R =
40\%$, $\rho = 0.30$, 100 names.  A 2{,}000-path Monte Carlo gives
a 5Y expected loss fraction $\approx 2.96\%$, matching the exact
$\mathrm{PD}(1-R) = 3.00\%$ within sampling error, independent of
$\rho$.  Unconditional expected loss is $\rho$-invariant; what
$\rho$ governs is the \emph{distribution} around that mean, hence
tranche prices.

Higher $\rho$ thickens both tails of the loss distribution:
simultaneous defaults and simultaneous survivals become more likely,
\emph{raising} senior-tranche expected loss (tail hits) and
\emph{lowering} equity-tranche expected loss (common survivals).
This is the core comparative static.

\begin{intuitionbox}
100 names, $\mathrm{PD} = 5\%$.  At $\rho = 0$: defaults cluster
near 5 with sd $\approx 2.2$, so a 10--15\% tranche is almost never
hit.  At $\rho = 1$: either all 100 default (5\%) or none do (95\%);
the super-senior 10--15\% tranche takes a full loss 5\% of the time.
Correlation pushes mass from the middle of the loss distribution
into both tails.
\end{intuitionbox}

\subsection{2008: the model that killed Wall Street}

The Gaussian copula was adopted as the industry standard right
after Li's 2000 paper because it was tractable: spreadsheet-ready
for rating agencies, mark-to-market-ready for dealers.  The CDO
market grew from nothing to \$4.7tn notional by 2007
\citep{salmon2009formula}.  When the US housing market cracked in
2007--08, the model's flaws were exposed catastrophically.

\paragraph{Failure 1: $\rho$ is regime-dependent.}  The copula
assumes a single constant $\rho$.  In reality, asset correlations
rise sharply in crises: 20\% in calm times becomes 70\% under
stress (``correlations go to one'').  Super-senior tranches priced
at low $\rho$ were massively mispriced when correlation spiked and
losses blew through the 15\% attachment point.

\paragraph{Failure 2: wrong calibration data.}  $\rho$ was
calibrated to 5Y equity-return correlations or index tranche
quotes, both predominantly normal-regime data.  Tail correlation
was systematically understated.

\paragraph{Failure 3: ratings depended on the model.}  Structured
product AAA ratings were themselves outputs of
Gaussian-copula-like models.  When the model failed, the rating
failed with it; AIG's super-senior CDS book required a \$180bn US
rescue in 2008 \citeHull.

\begin{historybox}[title={The formula that killed Wall Street}]
David X.\ Li's 2000 JP Morgan paper \emph{On Default Correlation:
A Copula Function Approach} brought the Gaussian copula to credit.
Within five years every rating agency, structured-product desk, and
risk system used it or a variant.  Felix Salmon's 2009 \emph{Wired}
piece ``Recipe for Disaster: The Formula That Killed Wall Street''
\citep{salmon2009formula} named it after the 2008 collapse of the
structured-credit market.  The paper is one of the most-cited and
most-vilified derivative-pricing notes in history.
\end{historybox}

\begin{pitfallbox}[The Gaussian copula is wrong in crisis]
The single-factor Gaussian copula prices tranches with a constant
$\rho$ from a static snapshot.  In a crisis:
\begin{itemize}
  \item asset correlations rise ($\rho \to 1$): tails are far
    heavier than the model assumes;
  \item the Gaussian has no tail dependence: extreme co-movements
    are near-zero-probability by construction;
  \item implied correlation across tranches becomes a surface, not
    a constant (the credit analogue of the vol smile).
\end{itemize}
Practitioners now use base-correlation models, $t$-copulas with
tail dependence, and stressed add-ons.  The Gaussian copula
survives as a \emph{quoting convention}, not a valuation truth.
\end{pitfallbox}


%% ================================================================
\section{Prosperity and desk connections}
\label{sec:credit_practical}
%% ================================================================

\subsection{CVA: credit risk in every derivative}
\label{sec:cva_preview}

Every OTC derivative has \emph{bilateral counterparty risk}: if
your swap is \$5M in the money and the counterparty defaults, you
become an unsecured creditor recovering at most \$5M $\times R$.
In plain English, \emph{CVA} is the market value of that
counterparty-default risk: the amount of value you would lose
today if the counterparty defaulted right now, weighted by how
likely they are to default and how far in the money the trade will
be when they do.

\begin{definition}[Credit Valuation Adjustment]
\index{CVA (credit valuation adjustment)|textbf}\index{CVA|textbf}
\emph{CVA} is the expected discounted loss on a derivative from
counterparty default.  For a single swap against counterparty $C$:
\[
  \mathrm{CVA}
  \;=\; (1 - R_C) \int_0^T D(t)\, \E_\Qm[\mathrm{EE}(t)]\, h_C(t)\, e^{-\int_0^t h_C(u)\, du}\, dt,
\]
with $\mathrm{EE}(t) = \max(V(t), 0)$ the expected positive exposure
(loss only when in our favour) and $h_C$ the counterparty hazard.
\end{definition}

CVA is structurally a CDS on your own derivative exposure, using
the same machinery (stripped hazards, recovery, $\int D\,d[1-S]$
integrals) layered on Monte Carlo exposure paths.  Mechanics are
deferred to \cref{ch:risk_production}; the point is that every desk
implicitly runs a CDS pricer, because every desk contributes to the
firm's counterparty-risk book.

\begin{intuitionbox}
Sell a \$10M IRS to a BB-rated hedge fund and you are long the
swap and \emph{short} CDS protection on the fund, a CDS on the
expected positive exposure, not on fixed notional.  CVA is the
price of that implicit CDS, and a bank's CVA desk is the internal
CDS-hedging desk the firm trades against.
\end{intuitionbox}

The symmetric adjustment is \emph{DVA} (debit value adjustment): the
expected discounted gain to us from the possibility that \emph{we}
default on the counterparty, computed with the same integral but
using our own hazard $h_{\mathrm{self}}$ and the expected
\emph{negative} exposure $\max(-V(t), 0)$.  Bilateral CVA from our
perspective is then $\mathrm{CVA} - \mathrm{DVA}$.  DVA is
controversial because it credits a bank's P\&L when its own credit
deteriorates, a ``benefit from our own default risk'' that is
impossible to hedge without buying protection on yourself.


\subsection{Mark-to-market of a seasoned CDS}
\label{sec:mtm_seasoned}

A CDS is costless at inception but acquires MtM as spreads move.
Enter a 5Y protection \emph{buyer} at $s_0 = 180$ bp; three months
later the residual 4.75Y reprices at $s_1 = 250$ bp.  The MtM is
\citeHull
\[
  V_{\mathrm{MtM}}
  \;=\; (s_1 - s_0) \cdot A(T - 0.25) \cdot N,
\]
where $A(T - 0.25)$ is the residual risky annuity of a 4.75-year
CDS at today's hazard curve.  Using the worked-example numbers and
a residual annuity of roughly 4.00, the MtM on \$10M notional is
\[
  V_{\mathrm{MtM}}
  \;\approx\; (0.025 - 0.018) \cdot 4.00 \cdot \$10\mathrm{M}
  \;=\; \$280{,}000,
\]
a profit for the buyer (insurance now worth more).  If spreads
tighten to 120 bp, MtM is $-\$240{,}000$.

\paragraph{Why annuity-weighted.}  $(s_1 - s_0)A$ follows from
$\mathrm{PL}(s) = s \cdot A$: the buyer pays $s_0 A N$ per annum
while the fair price is $s_1 A N$, so the buyer is $AN$ better off
per unit of widening.  This is the long-bond-rally mechanism ($A$
is the credit analogue of duration) and why CS01 $= A \cdot N
\cdot 10^{-4}$.

\begin{intuitionbox}
A CDS is an annuity-weighted short position on the spread.  Every
bp of widening is worth $A \cdot N \cdot 10^{-4}$ to the buyer:
\$4{,}200/bp on a 5Y \$10M CDS with $A \approx 4.2$.  A credit
trader's P\&L report is
``$\sum_i \text{CS01}_i \cdot \Delta s_i$'' over 250 names.
\end{intuitionbox}


\begin{prosperitybox}[title={Credit-curve stripping mirrors the zero-curve bootstrap}]
Prosperity has no defaulting instruments, but CDS stripping is
structurally identical to the Treasury-strip bootstrap
(\cref{ch:fixed_income_foundations}): a sequence of 1-D root-finds
across maturities converting observed prices into calibrated latent
rates.  Bootstrap a Prosperity coupon-bond curve and you can
bootstrap a CDS curve; only the instrument differs.
\end{prosperitybox}

\begin{gsbox}[title={The credit Strat's job description}]
The credit-derivatives Strat owns four systems, in order of
centrality: (i) a nightly/intraday \emph{PD stripping pipeline}
turning broker CDS quotes into per-entity hazard curves; (ii) a
\emph{single-name CDS pricer} with upfront/coupon conversion and
full Greeks (CS01, JTD); (iii) a \emph{CDO/tranche pricer}, Monte
Carlo under a Gaussian copula for sanity, base-correlation or
$t$-copula overlay for production; (iv) a \emph{CVA engine}
(\cref{ch:risk_production}) pricing counterparty-default risk on
every trade.  Interviewers test (i)--(ii) deeply and (iii)--(iv)
conceptually.
\end{gsbox}


%% ================================================================
\section{Key facts to memorise}
\label{sec:credit_summary}
%% ================================================================

\begin{itemize}
  \item A \textbf{CDS} is default insurance: the buyer pays a
    running premium $s$ on notional; the seller pays $(1-R) \times$
    notional on a credit event.
  \item \textbf{PD and survival}: $S(T) = 1 - \mathrm{PD}(T)$.
    Risk-neutral PD (implied by CDS) exceeds physical PD (historical
    agency default frequency) by a credit risk premium.
  \item \textbf{Recovery}: $R = 40\%$ for senior unsecured is the
    market convention.  LGD $= 1 - R$.
  \item \textbf{CDS par spread}:
    $s^\star = (1-R) \int_0^T D(t)\, d[1-S(t)] \,/\,
      \sum_i \delta_i D(t_i) S(t_i)$: expected discounted loss
    over risky annuity.
  \item \textbf{Quick approximation}: $s^\star \approx h (1-R)$.
    A 300 bp CDS with $R = 40\%$ implies a hazard rate of about
    500 bp.
  \item \textbf{Hazard rate}: $h(t)\,dt =
    \Prob(\tau \in [t, t+dt] \mid \tau > t)$.  Survival is
    $S(T) = \exp\!\left(-\int_0^T h(u)\, du\right)$.
  \item \textbf{Stripping} is the credit analogue of zero-curve
    bootstrapping (\cref{ch:fixed_income_foundations}): piecewise-
    constant hazard rates solved maturity-by-maturity, one root-find
    per step.
  \item \textbf{CDX.NA.IG} (125 IG names, fixed coupon 100~bp),
    \textbf{CDX.NA.HY} (100 HY names, 500~bp), \textbf{iTraxx Main}
    (125 European IG) are the flagship indices.
  \item \textbf{Tranches} carve a pool into seniority layers:
    equity (0--3\%), mezz, senior, super-senior (15\%+).  First
    losses hit the equity tranche.
  \item \textbf{Li's single-factor Gaussian copula}:
    $X_i = \sqrt{\rho} M + \sqrt{1-\rho} Z_i$, default times via
    $\tau_i = F_i^{-1}(\CDFN(X_i))$.  One $\rho$ for the whole
    pool.
  \item \textbf{2008 failure}: regime-dependent $\rho$ rises in
    crisis, so AAA super-senior tranches priced at low $\rho$ were
    catastrophically mispriced when defaults clustered: ``the
    formula that killed Wall Street'' \citep{salmon2009formula}.
  \item \textbf{Credit Greeks}: CS01 (1~bp credit sensitivity),
    JTD (jump-to-default), recovery delta.
\end{itemize}
